% GENERATED FILE. Edit canonical lessons, then run npm run generate:books.
% canonical-source-hash: fnv1a64:0275ad9fbb8915c9
% canonical-lessons: RU-C192-gazeta, RU-C192-zhurnal, RU-C192-slovar, RU-C192-lineyka, RU-C192-papka

\chapter{Newspaper, Magazine, Dictionary, Ruler, Folder}
\label{ch:ru-newspaper-magazine-dictionary-ruler-folder}
\hlchaptermodality{192}

\begin{chapteropening}
\textbf{By the end of this chapter:} \emph{I can say a newspaper, a magazine, a dictionary, a ruler and a folder.}

Complete the last lesson of chapter 192: I can say a newspaper, a magazine, a dictionary, a ruler and a folder.
\end{chapteropening}

\section[\texorpdfstring{\ru{газета} (gazéta)}{газета (gazéta)}]{\ru{газета} (gazéta) — a newspaper}

\label{lesson:RU-C192-gazeta}



\begin{quote}
Before the new one: say the Russian for a charger; morning exercises, then the Russian for a laptop.
\end{quote}

\subsection*{You'll want to know: \ru{газета}}
\textbf{\ru{газета}} (\emph{gazéta}) — \textquotedblleft{}a newspaper\textquotedblright{}.

\begin{grammarlens}[title={What you've built}]
One more word for this chapter.
\end{grammarlens}

\subsection*{Your turn}
\begin{itemize}
\raggedright
  \item \emph{Say it:} \emph{gazéta}
  \item \emph{Say it:} \emph{gazéta}, once more
  \item \emph{Say it:} say \emph{gazéta}
  \item \emph{Recall:} say the Russian for a charger; morning exercises, then the Russian for a laptop, then say \emph{gazéta} again
\end{itemize}

\begin{tcolorbox}[breakable,colback=teal!4,colframe=teal!35!black,title={Before you move on}]
What does \ru{газета} mean? (\textquotedblleft{}A newspaper\textquotedblright{}.) Say it once more.
\end{tcolorbox}
\section[\texorpdfstring{\ru{журнал} (zhurnál)}{журнал (zhurnál)}]{\ru{журнал} (zhurnál) — a magazine}

\label{lesson:RU-C192-zhurnal}



\begin{quote}
Before the new one: say the Russian for a laptop, then the Russian for a newspaper.
\end{quote}

\subsection*{You'll want to know: \ru{журнал}}
\textbf{\ru{журнал}} (\emph{zhurnál}) — \textquotedblleft{}a magazine\textquotedblright{}.

\begin{grammarlens}[title={What you've built}]
One more word for this chapter.
\end{grammarlens}

\subsection*{Your turn}
\begin{itemize}
\raggedright
  \item \emph{Say it:} \emph{zhurnál}
  \item \emph{Say it:} \emph{zhurnál}, once more
  \item \emph{Say it:} say \emph{zhurnál}
  \item \emph{Recall:} say the Russian for a laptop, then the Russian for a newspaper, then say \emph{zhurnál} again
\end{itemize}

\begin{tcolorbox}[breakable,colback=teal!4,colframe=teal!35!black,title={Before you move on}]
What does \ru{журнал} mean? (\textquotedblleft{}A magazine\textquotedblright{}.) Say it once more.
\end{tcolorbox}
\section[\texorpdfstring{\ru{словарь} (slovár)}{словарь (slovár)}]{\ru{словарь} (slovár) — a dictionary}

\label{lesson:RU-C192-slovar}



\begin{quote}
Before the new one: say the Russian for a newspaper, then the Russian for a magazine.
\end{quote}

\subsection*{You'll want to know: \ru{словарь}}
\textbf{\ru{словарь}} (\emph{slovár}) — \textquotedblleft{}a dictionary\textquotedblright{}.

\begin{grammarlens}[title={What you've built}]
One more word for this chapter.
\end{grammarlens}

\subsection*{Your turn}
\begin{itemize}
\raggedright
  \item \emph{Say it:} \emph{slovár}
  \item \emph{Say it:} \emph{slovár}, once more
  \item \emph{Say it:} say \emph{slovár}
  \item \emph{Recall:} say the Russian for a newspaper, then the Russian for a magazine, then say \emph{slovár} again
\end{itemize}

\begin{tcolorbox}[breakable,colback=teal!4,colframe=teal!35!black,title={Before you move on}]
What does \ru{словарь} mean? (\textquotedblleft{}A dictionary\textquotedblright{}.) Say it once more.
\end{tcolorbox}
\section[\texorpdfstring{\ru{линейка} (linéyka)}{линейка (linéyka)}]{\ru{линейка} (linéyka) — a ruler}

\label{lesson:RU-C192-lineyka}



\begin{quote}
Before the new one: say the Russian for a magazine, then the Russian for a dictionary.
\end{quote}

\subsection*{You'll want to know: \ru{линейка}}
\textbf{\ru{линейка}} (\emph{linéyka}) — \textquotedblleft{}a ruler\textquotedblright{}.

\begin{grammarlens}[title={What you've built}]
One more word for this chapter.
\end{grammarlens}

\subsection*{Your turn}
\begin{itemize}
\raggedright
  \item \emph{Say it:} \emph{linéyka}
  \item \emph{Say it:} \emph{linéyka}, once more
  \item \emph{Say it:} say \emph{linéyka}
  \item \emph{Recall:} say the Russian for a magazine, then the Russian for a dictionary, then say \emph{linéyka} again
\end{itemize}

\begin{tcolorbox}[breakable,colback=teal!4,colframe=teal!35!black,title={Before you move on}]
What does \ru{линейка} mean? (\textquotedblleft{}A ruler\textquotedblright{}.) Say it once more.
\end{tcolorbox}
\section[\texorpdfstring{\ru{папка} (pápka)}{папка (pápka)}]{\ru{папка} (pápka) — a folder}

\label{lesson:RU-C192-papka}



\begin{quote}
Before the new one: say the Russian for a dictionary, then the Russian for a ruler.
\end{quote}

\subsection*{You'll want to know: \ru{папка}}
\textbf{\ru{папка}} (\emph{pápka}) — \textquotedblleft{}a folder\textquotedblright{}.

\begin{grammarlens}[title={What you've built}]
One more word for this chapter.
\end{grammarlens}

\subsection*{Your turn}
\begin{itemize}
\raggedright
  \item \emph{Say it:} \emph{pápka}
  \item \emph{Say it:} \emph{pápka}, once more
  \item \emph{Say it:} say \emph{pápka}
  \item \emph{Recall:} say the Russian for a dictionary, then the Russian for a ruler, then say \emph{pápka} again
\end{itemize}

\begin{tcolorbox}[breakable,colback=teal!4,colframe=teal!35!black,title={Before you move on}]
What does \ru{папка} mean? (\textquotedblleft{}A folder\textquotedblright{}.) Say it once more.
\end{tcolorbox}
